%!TEX encoding = UTF-8 Unicode
%!TEX program = pdflatex
\documentclass[12pt, a4paper]{extarticle}

\usepackage{cmap}
\usepackage[utf8]{inputenc}
\usepackage[T2A]{fontenc}
\usepackage{mathptmx}
\DeclareFontFamily{T2A}{ptm}{}
\DeclareFontShape{T2A}{ptm}{m}{n}{<->sub*cmr/m/n}{}
\DeclareFontShape{T2A}{ptm}{bx}{n}{<->sub*cmr/bx/n}{}
\DeclareFontShape{T2A}{ptm}{b}{n}{<->sub*cmr/bx/n}{}
\DeclareFontShape{T2A}{ptm}{m}{it}{<->sub*cmr/m/it}{}
\DeclareFontShape{T2A}{ptm}{bx}{it}{<->sub*cmr/bx/it}{}
\DeclareFontShape{T2A}{ptm}{m}{sl}{<->sub*cmr/m/sl}{}
\DeclareFontShape{T2A}{ptm}{b}{it}{<->sub*cmr/bx/it}{}
\usepackage[a4paper,
  left=2cm,
  right=2cm,
  top=2cm,
  bottom=2cm]{geometry}
\usepackage[russian]{babel}

%% ===== Абзацные отступы =====
\setlength{\parindent}{0.8cm}
\setlength{\parskip}{0.4cm}

%% ===== Цвета =====
\usepackage{color}
\definecolor{dkgreen}{rgb}{0,0.6,0}
\definecolor{gray}{rgb}{0.5,0.5,0.5}
\definecolor{mauve}{rgb}{0.58,0,0.82}

%% ===== Листинги =====
\usepackage{listingsutf8}
\lstset{
  frame=none,
  language=C++,
  aboveskip=3mm,
  belowskip=3mm,
  showstringspaces=false,
  columns=flexible,
  basicstyle={\small\ttfamily},
  numbers=none,
  numberstyle=\tiny\color{gray},
  keywordstyle=\color{blue},
  commentstyle=\color{dkgreen},
  stringstyle=\color{mauve},
  breaklines=true,
  breakatwhitespace=true,
  tabsize=4,
  extendedchars=true,
  inputencoding=utf8/latin1
}

%% ===== Математика =====
\usepackage{mathtools}
\usepackage{amsmath}
\usepackage{amssymb}
\usepackage{amsthm}

%% ===== Иллюстрации =====
\usepackage{graphicx}
\graphicspath{{images/}}
\usepackage{caption}
\usepackage{subcaption}

%% ===== Оглавление =====
\usepackage{tocloft}
\renewcommand{\contentsname}{}
\renewcommand{\cftsecleader}{\cftdotfill{\cftdotsep}}

%% ===== Первый абзац =====
\usepackage{indentfirst}

%% ===== Списки =====
\usepackage{enumitem}

%% ===== Таблицы =====
\usepackage{diagbox}
\usepackage{booktabs}
\usepackage{longtable}
\usepackage{tabularx}
\usepackage{multirow}
\usepackage{array}

%% ===== Библиография =====
\usepackage[numbers,sort&compress]{natbib}
\bibliographystyle{plain}

%% ===== Микротипографика =====
\usepackage[protrusion=true,expansion=false]{microtype}
\emergencystretch=3em

%% ===== Гиперссылки =====
\usepackage[hidelinks]{hyperref}

%% ===== Русские подписи =====
\addto\captionsrussian{
  \renewcommand{\figurename}{Рисунок}
  \renewcommand{\tablename}{Таблица}
  \renewcommand{\refname}{Список литературы}
  \renewcommand{\appendixname}{Приложение}
}

%% ===== Пользовательские команды =====
\newcommand{\activeset}{\mathrm{active}}
\newcommand{\popcount}{\mathrm{popcount}}
\newcommand{\eval}[2]{\langle #1,\, #2 \rangle}
\newcommand{\reduce}{\;\rightarrow\;}

\title{}
\author{}
\date{}

% ================================================================
\begin{document}

%% Титульный лист
\begin{titlepage}

    \begin{center}
        {МИНИСТЕРСТВО НАУКИ И ВЫСШЕГО ОБРАЗОВАНИЯ \\
        РОССИЙСКОЙ ФЕДЕРАЦИИ}
        \\
        Федеральное государственное автономное образовательное учреждение высшего образования
        \\
        {\bfseries «Национальный исследовательский Нижегородский государственный университет им. Н.И. Лобачевского»\\(ННГУ)}
    
        \vspace{1em}
  
        {\bfseries Институт информационных технологий, математики и механики \\
        Кафедра: «Высокопроизводительных вычислений и системного программирования»} \\
        
        \vspace{1em}

        Направление подготовки: «Прикладная математика и информатика» \\
        Магистерская программа: «Вычислительные методы и суперкомпьютерные технологии»
    
    \end{center}


    \begin{center}
    \end{center}

    \vspace{5em}

    \begin{center}
        \textbf{ОТЧЕТ} \\
        по производственной преддипломной практике \\
        на тему: \\
        \textbf{«Создание инструмента для эмуляции исполнения CUDA PTX ядер на центральном процессоре»}
    \end{center}

    \vspace{6em}


    \begin{flushright}
        {\bfseries Выполнил:} студент группы\\3824М1ПМвм\\Мусин А.Н \underline{\hspace{3cm}} \linebreak\linebreak\linebreak
        {\bfseries Научный руководитель:}\\ доцент каф. ВВСП, к.т.н.\\ Горшков А.В.\underline{\hspace{3cm}} 
    \end{flushright}


    \vspace{\fill}

    \begin{center}
        Нижний Новгород\\2026
    \end{center}

\end{titlepage}

\setcounter{page}{2}

%% Содержание
\tableofcontents
\newpage

%% Главы
\section{Введение}

\subsection{Актуальность и мотивация}

Графические ускорители NVIDIA GPU занимают центральное место в современных
высокопроизводительных вычислениях: от обучения крупных языковых моделей до
физического моделирования и научных расчётов. Программирование GPU на платформе
CUDA строится на модели SIMT (Single Instruction, Multiple Threads), в которой
группы из 32~потоков~--- варпы~--- исполняют одну и ту же инструкцию над
различными данными~\cite{Lindholm2008, Nickolls2008}. Ключевым промежуточным
представлением в данной экосистеме служит PTX (Parallel Thread Execution)~---
аппаратно-независимый ассемблер NVIDIA~\cite{PTXISA}, в который компилируется
CUDA~C++ и из которого генерируется машинный код для конкретного поколения GPU.

Официальные инструменты NVIDIA (Nsight, Compute Sanitizer) требуют наличия
физического GPU и предоставляют агрегированную статистику на уровне SM,
не обеспечивая детальных метрик на уровне отдельных инструкций PTX.
Программный эмулятор PTX~ISA позволяет устранить оба ограничения: функционировать
на стандартном CPU без GPU-оборудования и собирать произвольные метрики
уровня инструкций.

Настоящий отчёт по практике посвящён двум взаимосвязанным аспектам реализации
такого эмулятора:
\begin{enumerate}
  \item разработке системы анализа производительности для эмулятора CUDA PTX;
  \item разработке тестовых CUDA-ядер и оценке корректности и производительности эмулятора.
\end{enumerate}

\subsection{Цели и задачи}

\textbf{Цель работы}~--- разработать систему профилирования, интегрированную в
программный эмулятор PTX~ISA, и верифицировать корректность эмулятора на
разнообразных реальных CUDA-ядрах.

Для достижения указанной цели поставлены следующие \textbf{задачи}:

\begin{enumerate}
  \item Спроектировать архитектуру системы метрик с кодогенерацией по JSON-манифесту,
        обеспечивающей добавление новой метрики без изменения C++-кода.
  \item Реализовать четыре ключевые метрики производительности:
        эффективность дивергенции ветвлений (\texttt{branch\_efficiency}),
        конфликты банков разделяемой памяти (\texttt{bank\_conflicts}),
        коэффициент слияния транзакций глобальной памяти (\texttt{global\_coalescing})
        и статистику записей в регистровый файл (\texttt{write\_ops}).
  \item Разработать инструмент генерации HTML-отчёта с аннотацией строк
        исходного \texttt{.cu}-файла по данным \texttt{.loc}-директив PTX.
  \item Разработать набор из семи интеграционных тестов на реальных CUDA-ядрах,
        перекрывающих качественно различные вычислительные паттерны.
  \item Провести верификацию корректности эмулятора и оценить производительность
        в сравнении с аппаратным GPU V100-PCIE.
\end{enumerate}

\subsection{Практическая значимость}

Разработанная система профилирования имеет непосредственное практическое применение:

\begin{itemize}
  \item \textbf{Детальное профилирование без GPU.} HTML-отчёт с аннотацией
        строк исходного кода позволяет определить «горячие» инструкции
        и оптимизировать шаблоны обращения к памяти на стандартном CPU.
  \item \textbf{Внедрение в конвейеры CI.} Эмулятор запускается как обычный
        Linux-процесс; для включения в \texttt{GitHub Actions} или
        \texttt{Jenkins} достаточно однострочного скрипта \texttt{cuemu ./binary}.
  \item \textbf{Образовательная ценность.} Метрики дивергенции и конфликтов
        банков обеспечивают наглядную демонстрацию механизмов SIMT-исполнения.
\end{itemize}

\newpage

\include{chapters/ch_profiling}
\include{chapters/ch_testing}
\section{Заключение}

В рамках практики решены две взаимосвязанные задачи: разработка системы
анализа производительности для эмулятора CUDA PTX и верификация
корректности эмулятора на реальных CUDA-ядрах.

\textbf{Система профилирования.} Реализован профилировщик с четырьмя
метриками (\texttt{branch\_efficiency}, \texttt{bank\_conflicts},
\texttt{global\_coalescing}, \texttt{write\_ops}), интегрированный в
эмулятор через кодогенерацию по JSON-манифесту. Архитектурное решение~---
хранение метаданных метрик в \texttt{profiling\_metrics.json} с автогенерацией
C++-кода скриптом \texttt{autogen.py}~--- позволяет добавлять новые метрики
без модификации горячего пути исполнения. Буферизация записей
\texttt{ProfilingRecord} на уровне варпа с единственным сбросом в общий
профилировщик по завершении блока минимизирует накладные расходы на синхронизацию.
Инструмент генерации HTML-отчёта (\texttt{renderer.py}) формирует
самодостаточный файл с пятью секциями, обеспечивая переход от сводных
метрик к аннотированному исходному коду в рамках одного документа.

\textbf{Верификация корректности.} Все семь интеграционных тестов пройдены
при максимальном отклонении от CPU-эталона не более $10^{-3}$.
Охваченные вычислительные паттерны: поэлементные операции (\texttt{vadd}),
трансцендентные функции (\texttt{gelu}), многофазные редукции (\texttt{softmax}),
тайловое матричное умножение (\texttt{mmul}), Flash Attention-подобное ядро,
FFT и итерационный метод сопряжённых градиентов. Верификация системы
профилирования подтвердила нулевые отклонения по всем шести аналитическим
тестам метрик.

\textbf{Производительность.} Замедление относительно V100-PCIE составляет
1\,900--2\,300$\times$ для нетривиальных нагрузок~--- неизбежная стоимость
программной эмуляции без JIT-компиляции. Накладные расходы профилировщика
при включённом сборе данных составляют 88--130\% для лёгких ядер и
500--800\% для вычислительно плотных (тайловое матричное умножение),
что приемлемо для задач анализа и отладки.

Среди текущих ограничений~--- неполнота покрытия PTX~ISA: ряд инструкций
не реализован вовсе (тензорные операции \texttt{mma} и \texttt{ldmatrix}),
а у реализованных поддержаны не все варианты модификаторов. Эмулятор
прошёл проверку на учебных примерах и лабораторных работах, однако
остаётся первой версией инструмента и требует существенного расширения
покрытия для применения на произвольных промышленных кодах.

\newpage


%% Список литературы
\nocite{Lindholm2008,Nickolls2008,VoltaWhitepaper,Graham1982}
\bibliography{references}

\label{lastpage}
\end{document}
